\documentclass[11pt]{article}

\usepackage[T1]{fontenc}
\usepackage{lmodern}
\usepackage[margin=1in]{geometry}
\usepackage{amsmath,amssymb,mathtools,amsthm}
\usepackage{booktabs}
\usepackage{microtype}
\usepackage[hidelinks]{hyperref}
\usepackage{url}

\allowdisplaybreaks[2]

\newtheorem{theorem}{Theorem}[section]
\newtheorem{proposition}[theorem]{Proposition}
\newtheorem{lemma}[theorem]{Lemma}
\newtheorem{corollary}[theorem]{Corollary}
\newtheorem{impinput}[theorem]{Imported Input}

\DeclareMathOperator{\tr}{tr}
\newcommand{\F}{\mathrm F}
\newcommand{\R}{\mathbb R}

\title{More than $67.3195\%$ of the zeros of the Riemann zeta function
are simple and lie on the critical line}
\author{Nicholas Pipitone\\
\href{mailto:npip99@gmail.com}{\texttt{npip99@gmail.com}}\\
\url{https://github.com/npip99}}
\date{August 11, 2026}
\hypersetup{
  pdfauthor={Nicholas Pipitone},
  pdftitle={More than 67.3195 percent of the zeros of the Riemann zeta function
  are simple and lie on the critical line}
}

\begin{document}

\maketitle

\begin{abstract}
Let $N(T,2T)$ count the zeros of the Riemann zeta function with ordinates in
$(T,2T]$, with multiplicity, and let $N_0^s(T,2T)$ count the simple zeros on
the critical line in that interval.  We prove
\[
 \liminf_{T\to\infty}\frac{N_0^s(T,2T)}{N(T,2T)}
 \ge 0.673195198901\ldots \;>\; \frac{673195}{10^6}.
\]
The proof combines the analytic framework of Claude~\cite{Claude26}, the
stability refinement of the rank--trace argument
introduced in~\cite{Ainta26}, and the re-optimized window, position-weighted
seven-point inequality, and sharp block profile of~\cite{Trmdy26}.  The single
new mathematical ingredient is a strengthened certified constant for the
weighted seven-point inequality of~\cite{Trmdy26}: we certify
$F\ge 509/100000$ where $1/200$ was certified before, and re-optimize the
block length.  The finite claims are certified by rigorous interval
computations, with the principal certificate repeated at two grid
resolutions; the deduction is reproduced in full.
\end{abstract}

\section{Statement and imported inputs}\label{sec:setup}

Throughout, write $N:=N(T,2T)$ for the number of zeros of $\zeta$ with
ordinates in $(T,2T]$, counted with multiplicity, and
$S:=N_0^s(T,2T)$ for the simple zeros on the critical line in that range.

\begin{theorem}\label{thm:main}
\[
 \liminf_{T\to\infty}\frac{N_0^s(T,2T)}{N(T,2T)}
 \;\ge\;
 \frac{250\,H_{\mathrm{cert}}-\eta\,\frac{3}{1150}\cdot 249}{250-R}
 \;=\;0.6731951989015205\ldots\;>\;\frac{673195}{10^{6}},
\]
where $H_{\mathrm{cert}}=672457/10^6$,
$A=31049/25000$, $R=2\sqrt A-1$, and $\eta=R/A$.
\end{theorem}

The proof template is that of~\cite{Trmdy26}, which in turn builds
on~\cite{Ainta26} and~\cite{Claude26}.  We isolate below the two external
analytic consequences used in the finite deduction.  The first is elementary
from the Riemann--von Mangoldt formula; the second combines the arbitrary-
window trace asymptotics of~\cite[\S7.1]{Claude26} with the stability
rank--trace refinement of~\cite{Ainta26}.

\begin{impinput}[Normalized ordinates and total span]\label{inp:span}
Put $\ell_T=\log(T/(2\pi))$ and, for a zero of ordinate $\gamma$, put
\[
 x_\gamma=\frac{\ell_T(\gamma-T)}{2\pi}.
\]
After the deletion of $o(N)$ exceptional zeros, the retained simple zeros
have normalized ordinates $x_1<\cdots<x_{S^\circ}$ with
$S^\circ=S-o(N)$ and
\begin{equation}\label{eq:totalspan}
 x_{S^\circ}-x_1\le\frac{T\ell_T}{2\pi}=N+o(N).
\end{equation}
Here the second relation is asymptotic.  In the present notation the
Riemann--von Mangoldt formula gives
\[
 N(T,2T)=\frac{T}{2\pi}\bigl(\ell_T+2\log2-1\bigr)+O(\log T)
 =\frac{T\ell_T}{2\pi}+O(T)
 =\frac{T\ell_T}{2\pi}+o(N).
\]
\end{impinput}

\begin{impinput}[Stability counting inequality]\label{inp:stability}
Let $v:[-\tfrac12,\tfrac12]\to[\tfrac34,1]$ be an even $C^2$ function that
is nonincreasing on $[0,\tfrac12]$.  At the endpoint $\lambda=1$, put
$\mathcal L=\ell_T$, fix a ramp width $w_0\ge1$ independently of $T$, and use
$\sqrt{v(u/\mathcal L)}$ in the test family of~\cite[\S2]{Claude26},
multiplied by its standard even boundary ramp of support
$[-\mathcal L/2,\mathcal L/2]$.  The required condition
$8w_0\le\mathcal L$ then holds for all sufficiently large $T$.  This is an
admissible window provided the variation and derivative conditions recorded
in~\cite[\S7.1]{Claude26} hold; the endpoint $\lambda=1$ is explicitly
allowed by~\cite[Remark~6.1]{Claude26}.  Write
\[
 K_v(x)=\int_{-1/2}^{1/2}v(t)\cos(2\pi xt)\,dt,\qquad
 k_v=\frac{K_v}{K_v(0)},\qquad
 H(v)=2-\frac1{c_1(v)},
\]
where $c_1(v)=\bigl(\int v\bigr)^2/\bigl(\int v^2+\iint|s-s'|v(s)v(s')\,ds\,ds'\bigr)$.
Then
\begin{equation}\label{eq:stability}
 S\;\ge\;H(v)\,N+\tr\Psi(M)-o(N),
\end{equation}
where $M$ is the Gram matrix of the retained simple-zero atoms,
$\Psi(t)=(t-1)^2$ for $0\le t\le2$ and $2t-3$ for $t\ge2$, and the entries of
$M$ satisfy $M_{ij}=k_v(x_i-x_j)+o(1)$ uniformly over all retained pairs in
the dyadic interval whose normalized ordinate differences lie in any fixed
compact set, the exceptional $o(N)$ zeros being deleted.
\end{impinput}

We use one retained set for both imported inputs, enlarging the exceptional
set to the union of the two $o(N)$ sets if necessary; it still has size
$o(N)$.  This further deletion replaces the original Gram matrix by a
principal compression.  Pinching into retained and discarded principal
blocks, followed by $\Psi\ge0$, gives
\[
 \tr\Psi(M_{\rm original})
 \ge\tr\Psi(M_{\rm retained})+\tr\Psi(M_{\rm discarded})
 \ge\tr\Psi(M_{\rm retained}),
\]
so the stability inequality remains valid with the smaller retained matrix.
Further deletion preserves the order of the ordinates and can only decrease
their outer span; thus Input~\ref{inp:span} also remains valid for this common
retained set, whose consecutive gaps are used below.

Input~\ref{inp:stability} is the stability form of the rank--trace argument:
the analytic content (explicit formula, Gabor trace asymptotics for
admissible windows, tail estimates, and Gram-entry asymptotics) is from
\cite{Claude26}; retention of the defect $\tr\Psi(M)$ is the contribution
of~\cite{Ainta26}, with the combined interface also used in~\cite{Trmdy26}.
The deduction below uses \eqref{eq:stability} only through the block/pinching
mechanism of Section~\ref{sec:deduction}.

\section{The window}\label{sec:window}

We use the window of~\cite{Trmdy26} unchanged:
\begin{equation}\label{eq:window}
 v(s)=\sum_{j=0}^{6}c_j\cos(\omega_j s),\qquad
 \omega=(\sqrt2,\,2\pi,\,4\pi,\,6\pi,\,8\pi,\,10\pi,\,12\pi),
\end{equation}
\[
 c=\bigl(10^9,\;3322500,\;-7609135,\;1190194,\;-731476,\;-1680572,\;1141360\bigr)/10^9 .
\]

\begin{proposition}[Certified window bounds]\label{prop:window}
On $[-\tfrac12,\tfrac12]$ one has $\tfrac34\le v\le1$
(certified global bounds: $v\ge0.750213\ldots$ and
$v\le0.995632\ldots$), and
on $[0,\tfrac12]$ one has $v'\le0$.  Moreover,
\[
 H(v)=0.67245704141454\ldots\;\ge\;H_{\mathrm{cert}}:=\frac{672457}{10^6}.
\]
\end{proposition}

\begin{proof}[Verification]
Arb interval arithmetic over the exact defining data.  The range enclosure
uses $8192$ cells.  For monotonicity, the verifier proves $v''<0$ on the
contiguous union of the four cells adjacent to zero; since the cosine sum is
even, $v'(0)=0$, so $v'<0$ away from the origin on that union.  It proves
$v'<0$ directly on every remaining cell.  The enclosures are reproduced by
the \texttt{fast} module of the verifier (see Section~\ref{sec:repro}).
The bounds for $v$ and $H(v)$ reproduce the corresponding certificate
of~\cite{Trmdy26}; the derivative check makes explicit the admissibility
condition used here.
\end{proof}

We now verify the applicability clause in Input~\ref{inp:stability}.  Let
\(\rho\in C^3(\mathbb R)\) be the fixed nondecreasing function of
\cite[\S2.2]{Claude26}, equal to zero on $(-\infty,0]$ and to one on
$[1,\infty)$.  Define the smooth even boundary ramp by
\[
 \rho_{\mathcal L,w_0}(u)
 =\rho\!\left(\frac{\mathcal L/2-|u|}{w_0}\right)
\]
and put
\[
 \varphi_{\mathcal L,w_0}(u)=
 \begin{cases}
  \sqrt{v(u/\mathcal L)}\,\rho_{\mathcal L,w_0}(u),
    & |u|\le\mathcal L/2,\\
  0,&|u|>\mathcal L/2.
 \end{cases}
\]
Since $v\ge3/4$, its square root is smooth.  Proposition~\ref{prop:window}
and the defining properties of the ramp show that both
$\varphi_{\mathcal L,w_0}$ and $\varphi_{\mathcal L,w_0}^2$ are even,
nonnegative, compactly supported, $C^2$, bounded
by one, and nonincreasing for $u\ge0$.  Their total variations are therefore
$2\varphi_{\mathcal L,w_0}(0)\le2$ and
$2\varphi_{\mathcal L,w_0}(0)^2\le2$.  On the fixed-width
boundary strips, the product rule and the ramp estimates give
$\|\varphi_{\mathcal L,w_0}''\|_1+
\| (\varphi_{\mathcal L,w_0}^2)''\|_1=O(1/w_0)$; in the interior,
scaling the fixed smooth functions $\sqrt v$ and $v$ contributes only
$O(1/\mathcal L)$.  Thus, for fixed $w_0$ and all sufficiently large
$\mathcal L$, this family
satisfies the support, variation, and second-derivative hypotheses
of~\cite[\S7.1]{Claude26}.  This closes the arbitrary-window applicability
step rather than treating it as an additional numerical assumption.  Finally,
the boundary strips have total length $O(w_0)$, while the full support has
length $\mathcal L$.  Their contribution to the normalized moments in~(7.1)
of~\cite{Claude26} is therefore $O(w_0/\mathcal L)=o(1)$; likewise the
normalized autocorrelation
kernel converges, uniformly for bounded normalized shifts, to $k_v$ by the
$L^1$ bound on the ramp discrepancy.  Thus the limiting constant and kernel
in Input~\ref{inp:stability} are precisely $H(v)$ and $k_v$ displayed above.

\section{The sharp block profile}\label{sec:profile}

\begin{lemma}\label{lem:profile}
Let $G\succeq0$ be Hermitian with $E:=2\sum_{i<j}|G_{ij}|^2$.  Then
\[
 \tr\Psi(G)\;\ge\;h(E),\qquad
 h(E)=\begin{cases}E,&0\le E\le1,\\ 2\sqrt E-1,&E\ge1,\end{cases}
\]
and $h$ is sharp.
\end{lemma}

\begin{proof}
For $\lambda\ge0$ write $x=(\lambda-1)^2$.  Since
$\Psi(\lambda)=(\lambda-1)^2-\bigl((\lambda-2)_+\bigr)^2$ and, for
$\lambda\ge1$, $(\lambda-2)_+=(\sqrt x-1)_+$ while for $0\le\lambda<1$ both
$(\lambda-2)_+=0$ and $x\le1$, one has the exact identity
$\Psi(\lambda)=g\bigl((\lambda-1)^2\bigr)$ with
\[
 g(x)=\begin{cases}x,&0\le x\le1,\\ 2\sqrt x-1,&x\ge1.\end{cases}
\]
The function $g$ is nonnegative, increasing, concave, and $g(0)=0$; hence $g$
is subadditive, so for the eigenvalues $\lambda_i$ of $G$,
\[
 \tr\Psi(G)=\sum_i g\bigl((\lambda_i-1)^2\bigr)
 \;\ge\;g\Bigl(\sum_i(\lambda_i-1)^2\Bigr)
 =g\bigl(\|G-I\|_{\F}^2\bigr)\;\ge\;g(E),
\]
using that $\|G-I\|_{\F}^2=\sum_i(\lambda_i-1)^2$ and that the off-diagonal
part alone contributes $E$.  For sharpness, put
$t_n=\sqrt{E(n-1)/n}$, $a_n=1+t_n$, and
$b_n=1-t_n/(n-1)$, taking $n$ large enough that $b_n\ge0$.  With
$u_n=n^{-1/2}(1,\ldots,1)$, set
$G_n=b_nI+(a_n-b_n)u_nu_n^*$.  This matrix is positive semidefinite with
constant diagonal one and off-diagonal energy exactly $E$.  Its eigenvalues
are $a_n,b_n,\ldots,b_n$, and therefore
\[
 \tr\Psi(G_n)=\Psi(1+t_n)+(n-1)\Psi\bigl(1-t_n/(n-1)\bigr)
 \longrightarrow h(E).
\]
Thus no larger lower profile is possible even among correlation matrices.
\end{proof}

\begin{corollary}[Chord bound]\label{cor:chord}
For $0\le E\le A$ one has $h(E)\ge\frac{h(A)}{A}\,E$; and $h$ is
nondecreasing.
\end{corollary}

\begin{proof}
The first assertion is the chord inequality for the concave function $h$
between $(0,0)$ and $(A,h(A))$; monotonicity is immediate from its displayed
formula.
\end{proof}

\section{The certified weighted seven-point inequality}\label{sec:seven}

Let $w=k_v^2$ for the window \eqref{eq:window}, let $g_1,\dots,g_6\ge0$ be
consecutive normalized gaps with partial sums $y_j=g_1+\dots+g_j$
($y_0=0$), and define
\begin{equation}\label{eq:F}
 F(g_1,\dots,g_6)
 =\frac{1}{2300}\sum_{i=1}^{6}g_i
 +\sum_{0\le i<j\le6}a_{ij}\,w(y_j-y_i),
\end{equation}
with the reflection-symmetric exact rational weights $a_{ij}$ of
Table~\ref{tab:weights}.  Every span capacity is exactly
$\sum_{i}a_{i,i+r}=2$ ($r=1,\dots,6$).

\begin{table}[h]
\centering
\begin{tabular}{ccccccc}
\toprule
$10^6a_{ij}$ & $j=1$ & $2$ & $3$ & $4$ & $5$ & $6$\\
\midrule
$i=0$ & 239252 & 528172 & 965879 & 1000000 & 1000000 & 2000000\\
$1$ & & 381335 & 465776 & 34121 & 0 & 1000000\\
$2$ & & & 379413 & 12104 & 34121 & 1000000\\
$3$ & & & & 379413 & 465776 & 965879\\
$4$ & & & & & 381335 & 528172\\
$5$ & & & & & & 239252\\
\bottomrule
\end{tabular}
\caption{The pair weights of \cite{Trmdy26}, used unchanged.}
\label{tab:weights}
\end{table}

\begin{proposition}[Certified local inequality]\label{prop:F}
For all $g_1,\dots,g_6\ge0$,
\begin{equation}\label{eq:Fbound}
 F(g_1,\dots,g_6)\;\ge\;\varepsilon_0:=\frac{509}{100000}=0.00509 .
\end{equation}
\end{proposition}

\begin{proof}[Verification]
The pressure term first reduces the unbounded orthant to a compact simplex:
if $\sum_i g_i\ge\varepsilon_0/p=11707/1000$, then the pressure alone proves
\eqref{eq:Fbound}.  On the complementary simplex, use exhaustive interval
subdivision with Arb enclosures of $w$ on a grid of mesh
$1/4000$, exactly the decision procedure of~\cite{Trmdy26} with the target
raised from $1/200$ to $509/100000$: one-body pruning, per-box range-minimum
lower bounds over all $21$ pair separations, a certified convex tangent
bound, and bisection of the widest coordinate; the run fails loudly at any
unresolved terminal cell.  The recorded run has $1{,}739{,}356$ nodes and maximum depth $57$; the
weaker target $127/25000$ was additionally certified in separate regression
runs.  See Section~\ref{sec:repro} for the recorded
certificates.  (The observed minimum of $F$ is $\approx0.005091029$, at a
lattice-like configuration.)
\end{proof}

\section{Deduction of the main theorem}\label{sec:deduction}

Set
\[
 q=6,\qquad m=250,\qquad p=\frac1{2300},\qquad B_p=6p=\frac{3}{1150},
\]
\[
 A=\varepsilon_0\,(m-q)=\frac{509}{100000}\cdot244=\frac{31049}{25000}
 =1.24196,\qquad
 R=h(A)=2\sqrt A-1,\qquad \eta=\frac RA .
\]
Numerically $R=1.22886518210501\ldots$, $\eta=0.98945632879079\ldots$.

\subsection{Windows to blocks}

Let $B$ be a block of $m$ consecutive retained simple zeros with normalized
ordinates $y_1<\dots<y_m$, Gram matrix $G_B$, and
$E_B=2\sum_{i<j}|(G_B)_{ij}|^2$.  Summing \eqref{eq:Fbound} over the $m-q$
consecutive seven-point windows of $B$: a pair spanning $r$ gaps lies in at
most $7-r$ windows and its weight totals at most the span capacity $2$; each
gap lies in at most $6$ windows, contributing at most
$6p\sum g_i=B_p\,\mathrm{span}(B)$.  By the nonnegativity of the weights,
this gives the purely kernel-level bound
\begin{equation}\label{eq:kernelblock}
 2\sum_{i<j}w(y_j-y_i)+B_p\,\mathrm{span}(B)\;\ge\;A.
\end{equation}

\subsection{Blocks to defect}

We claim
\begin{equation}\label{eq:blockdefect}
 \tr\Psi(G_B)+\eta\,B_p\,\mathrm{span}(B)\;\ge\;R-o(1).
\end{equation}
Put
\[
 D:=A/B_p=714127/1500.
\]
If $\mathrm{span}(B)\ge D$, \eqref{eq:blockdefect} is immediate:
$\tr\Psi(G_B)\ge0$ and
$\eta B_p\,\mathrm{span}(B)\ge\eta A=R$.  We may therefore suppose that
$\mathrm{span}(B)<D$.  Every pair separation in $B$ then lies in the fixed
compact interval $[0,D]$.  The Gram-entry asymptotic of
Input~\ref{inp:stability} is uniform there.  More explicitly, its error is
bounded by some $\delta_T(D)\to0$ for every pair in every such block.  Since
$m$ is fixed, $|k_v|\le1$ because $v\ge0$, and the same asymptotic gives
$|(G_B)_{ij}|\le2$ for all sufficiently large $T$.  Consequently
\[
 2\sum_{i<j}w(y_j-y_i)
 \le E_B+O(m^2\delta_T(D))=E_B+o(1).
\]
Together with \eqref{eq:kernelblock}, this gives
\begin{equation}\label{eq:blockenergy}
 E_B+B_p\,\mathrm{span}(B)\;\ge\;A-o(1).
\end{equation}
If $E_B\ge A$, then Lemma~\ref{lem:profile} and monotonicity give
$\tr\Psi(G_B)\ge h(E_B)\ge h(A)=R$.  If $E_B<A$, the chord bound
(Corollary~\ref{cor:chord}) gives
$\tr\Psi(G_B)\ge\eta E_B$, and \eqref{eq:blockenergy} yields
$\eta E_B+\eta B_p\,\mathrm{span}(B)\ge\eta A-o(1)=R-o(1)$.
Because $m$ and $D$ are fixed, the $o(1)$ here is uniform over the
small-span blocks; the large-span case has no asymptotic error.

\subsection{Pinching and averaging}

Exactly as in \cite[\S5]{Ainta26} (and \cite[\S4]{Trmdy26}): for each of the
$m$ offsets, partition the retained simple zeros into consecutive full
blocks of size $m$, leaving at most $2(m-1)$ endpoint points.  If
$\mathcal P_rM$ denotes the resulting block-diagonal pinching (including the
remaining endpoint block), then pinching is an average of unitary
conjugations.  Convexity, unitary invariance, and $\Psi\ge0$ therefore give
the precise inequality
\[
 \tr\Psi(M)\ge\tr\Psi(\mathcal P_rM)
 \ge\sum_{B\text{ full at offset }r}\tr\Psi(G_B).
\]
The average number of full blocks over the offsets is
$S^\circ/m+O(1)=S/m+o(N)$.  There are $O(N)$ blocks because $S\le N$, so the
uniform per-block $o(1)$ above sums to $o(N)$.  Each interior gap lies in a
block span for at most $m-1$ offsets, and Input~\ref{inp:span} bounds their
total length by $N+o(N)$.  Summing \eqref{eq:blockdefect} over full blocks
and averaging over the offsets therefore gives
\begin{equation}\label{eq:defect}
 \tr\Psi(M)\;\ge\;\frac Rm\,S-\eta\,B_p\,\frac{m-1}{m}\,N-o(N).
\end{equation}

\subsection{Conclusion}

Insert \eqref{eq:defect} and Proposition~\ref{prop:window} into
\eqref{eq:stability}:
\[
 S\;\ge\;H_{\mathrm{cert}}\,N+\frac Rm S-\eta B_p\frac{m-1}{m}N-o(N),
\]
hence
\[
 \liminf_{T\to\infty}\frac SN
 \;\ge\;
 \frac{m\,H_{\mathrm{cert}}-\eta\,B_p\,(m-1)}{m-R}
 =\frac{250\cdot\frac{672457}{10^6}-\eta\cdot\frac{3}{1150}\cdot249}{250-R}
 =0.6731951989015205755\ldots
\]
The final arithmetic (including the enclosure of $R$ and $\eta$ and the
comparison with $673195/10^6$) is itself certified in Arb by the verifier's
\texttt{fast} module.  This proves Theorem~\ref{thm:main}. \qed

If $\mathcal N(X)$ counts zeros up to height $X$ with multiplicity and
$\mathcal S(X)$ counts the simple critical-line zeros up to height $X$, the
same constant also satisfies
\[
 \liminf_{X\to\infty}\frac{\mathcal S(X)}{\mathcal N(X)}
 \ge0.6731951989015205\ldots .
\]
Indeed, decompose $(0,X]$ into dyadic intervals
$(X/2^{j+1},X/2^j]$.  For any fixed error tolerance, Theorem~\ref{thm:main}
applies to every interval whose lower endpoint is sufficiently large; summing
their inequalities shows that the omitted fixed bottom interval is negligible
as $X\to\infty$, and then the error tolerance may tend to zero.  This is the
standard dyadic-summation passage from the
displayed theorem to the density formulation in the title.

\subsection{Sanity gates}

With $(\varepsilon_0,p,m)=(1/200,\,1/2300,\,257)$ the same formula gives
$0.673137630699\ldots$, the bound of~\cite{Trmdy26}; with the unit-cap
profile ($\eta=1$, $R=A\le1$) and the parameters of~\cite{Ainta26}
($\varepsilon_0=19/5000$, $p=1/3000$, $m=269$, Montgomery--Taylor window) it
gives $0.673008527927\ldots$.  These are regression checks on the deduction;
they are not additional inputs to the present bound.

\section{Reproducibility}\label{sec:repro}

All finite claims (Propositions~\ref{prop:window} and~\ref{prop:F} and the
final arithmetic) are certified by rigorous interval computation.  Arb
\cite{Johansson17} encloses the transcendental evaluations.  The range-
minimum path stores outward-rounded binary64 lower bounds derived from Arb
intervals and combines nonnegative terms with directed rounding; the tangent
path separately uses Arb enclosures for signed derivatives, gradients, and
the LDL convexity test.  Source,
recorded certificates, and instructions are at
\begin{center}
 \url{https://github.com/npip99/zeta-zeros}
\end{center}
These computations do not certify the external analytic statements isolated
as Inputs~\ref{inp:span} and~\ref{inp:stability}; those remain dependencies on
the cited sources, as is explicit in Section~\ref{sec:setup}.
The verifier is that of~\cite{Trmdy26}, vendored under MIT with the design
constants and an additional monotonicity check changed.  As a resolution
cross-check, the same decision procedure certified Proposition~\ref{prop:F}
twice: the headline target
$509/100000$ at grid $1/4000$ ($1{,}739{,}356$ nodes, depth 57) and again at
grid $1/8000$ ($1{,}736{,}620$ nodes, depth 57), the intermediate target
$127/25000$ also certified in two separate runs (including a
$1{,}544{,}746$-node run during the design exchange).
These repeated runs share an implementation and are not claimed to be
independent formal proofs.  The grid parameter controls the one-dimensional
kernel enclosure tables, not the initial six-dimensional box cover, so a
similar subdivision-tree size at the two resolutions is expected.

\subsection*{Credit}

This note is a minimal increment on~\cite{Trmdy26} (whose window, weights,
profile, and base verifier are reused), which in turn extends
\cite{Ainta26} over the framework of~\cite{Claude26}.  The improvement was
found and cross-verified in a Claude--GPT-5.6 collaboration on August 11,
2026.

\begin{thebibliography}{99}
\small

\bibitem{Claude26}
Claude,
\emph{More than two thirds of the zeros of the Riemann zeta function lie on
the critical line},
preprint, August 10, 2026,
\href{https://www-cdn.anthropic.com/564f962e60643842f5fcb4a17c9dbc8f608f1c37.pdf}{available online}.

\bibitem{Ainta26}
ainta,
\emph{zeta-simple-zeros: a 67.30085\% lower bound for simple zeros of the
Riemann zeta function},
research artifact, August 10, 2026,
\url{https://github.com/ainta/zeta-simple-zeros}.

\bibitem{Trmdy26}
T.~Haugland et al.,
\emph{zeta-simple-zeros-673137: a 67.3137630699\% lower bound for simple
zeros of the Riemann zeta function},
research artifact, August 11, 2026,
\url{https://github.com/trmdy/zeta-simple-zeros-673137}.

\bibitem{Johansson17}
F.~Johansson,
\emph{Arb: efficient arbitrary-precision midpoint-radius interval
arithmetic},
IEEE Trans.\ Comput.\ \textbf{66} (2017), no.~8, 1281--1292.

\end{thebibliography}

\end{document}
